\scanpage{286}
\mspage{202}
% archive sheet number: 285
Pour ceci, on utilise (4.7.6.2) page 171

\[ M_{\unsure{\pi|\kappa|k}}\subset\mathrm{Filt}^{\pi^{\pi_0}}\qquad(\text{si }\pi\geq\pi_0,\pi^{\pi_0}\geq\pi_1), \]

et la définition de $\mathrm{Filt}^{\pi'}$ (pour $\pi'\geq\pi_0$) \deleted{\unsure{donne}} montre que c'est là la condition cherchée.

Prouvons donc $2^\circ)$. Comme $M_{\pi_0}$ est stable par \inserted{les} $\varinjlim$ inductives finies qui sont représentables dans $M$ (SGA 4 I \unsure{9.9.}), il est stable par facteurs directs, et comme il contient $C$, on a donc $C^{\unsure{\natural|4}}\subset M_{\pi_0}$. \deleted{\unsure{Pour l'i}} L'inclusion inverse $M_{\pi_0}\subset C^{\natural}$ se voit de la façon \unsure{standard|standart}, en écrivant un objet \inserted{$F$} de $M_{\pi_0}$ comme $\varinjlim_I x_i$, les $x_i$ dans $C$, $I$ grd devant $\pi_0$, alors $F$ \deleted{\unsure{se}} est facteur direct d'un des \unsure{$x_i$|ci}.

Le fait que $C^{\natural}$ est une \uline{enveloppe de Karoubi} de $C$ dans $M\hookrightarrow C^{\wedge}$ provient du fait que $M$ est stable par facteurs directs, et que $C$ est une sous-catégorie pleine de $M$, cf SGA 4 I (?...).

4.8.6. \inserted{\uuline{\unsure{Questions ouvertes|Question ouverte}}} On a ainsi prouvé 4.8.1 revu et corrigé. Il reste les questions \deleted{de}

$1)$ \unsure{Donner} un exemple où \inserted{4.8.1} a) est \unsure{bien|elle} vérifiée, mais non b) etc.

% slanted
\margin{Il suffit de trouver un cardinal non régulier \deleted{\unsure{pour}} pour \unsure{la} feg \unsure{canon.} de $\mathrm{Ind}_{\pi_0}C$.}

$2)$ Dégager des conditions supplémentaires \inserted{\unsure{serviables|vérifiables|raisonnables}} qui, jointes à a), impliquent b), donc c, c'.
